\section{アプローチ10（ONNX Runtime WebによるYOLOv12 Nanoへの検出モデル刷新と推論パイプライン再構築）}

\subsection{手法の概要}
アプローチ09までの複数パイプライン（一次検出・拡大・ポーズ検出・各種重複排除）における複雑な計算ステップと処理オーバーヘッドを解消するため、アプローチ10では人物検出エンジンを完全に構築し直し、ONNX Runtime Web 上で動作する高精度かつ軽量な YOLOv12 Nano モデルへの移行を実施した。

検出精度の飛躍的向上に伴い、アプローチ03〜08で採用していた従来の段階的重複排除処理（\texttt{deduplicateDetections.ts}）を完全に廃止し、YOLOv12 Nano の出力特性に特化した高速な重複排除アルゴリズムへ刷新した。

本アプローチにおける検出・重複排除およびシステム構造の最適化仕様は以下の通りである。

\begin{itemize}
    \item \textbf{確信度フィルタリングとメモリ解放}：検出信頼度（例: 0.25）による事前フィルタリングを行い、推論過程で生成されるテンソルやバッファは \texttt{dispose} メソッド等を用いて確実に解放した。また、メモリ制限環境に対応するため、必要に応じてアリーナ割り当ての無制限化を無効に設定した。
    \item \textbf{座標ソートと横方向被りによるバウンディングボックス統合}：
    検出されたバウンディングボックス群をX座標順に並び替えた後、順次比較を行う。横方向の位置および幅が設定閾値（例: 8割）以上重なっている場合（※内包関係は直接の除外条件とはせず、閾値判定のロジックに従い処理）、それらを1人の人物領域として統合・マージする。マージ後のバウンディングボックスは、統合対象となった全領域を包含する縦・横の最大範囲として決定され、使用された元のバウンディングボックスは配列から削除される。この操作を反復し、グループ判定の前段となる人物領域を全確定させる。
    \item \textbf{DOM非依存化とWorker管理の全面移行}：MediaPipeとは異なり ONNX Runtime Web は DOM に依存しないため、人物検出処理自体を完全に Web Worker（\texttt{peopleDetectionWorker.ts}）へと移行させた。これに伴い、\texttt{workerManager.ts} が人物検出 Worker も一元管理する構造に再設計し、従来のポーズ検出 Worker における \texttt{PROCESS\_CANDIDATE} メッセージ仕様は廃止した。
    \item \textbf{モジュール構造の再編}：ディレクトリ構成を見直し、ルートの \texttt{detectPeople.ts} をバレルファイルから Worker 呼び出しを仲介する単一のエントリーポイントへと役割を変更・保持させた。
\end{itemize}

\subsection{結果およびアプローチ09との比較・課題}

\subsubsection{アプローチ09との比較}
YOLOv12 Nano への移行とX座標ソートベースの高速統合アルゴリズムの導入により、一次検出の精度と処理速度が大幅に向上した。また、DOM依存を排除して Worker スレッドへ全処理を移管したことで、メインUIスレッドの負荷がさらに軽減され、メモリ管理の堅牢性と描画パフォーマンスの両立が達成された。

\subsubsection{判明した課題}
\begin{enumerate}
    \item \textbf{待機列内における前後跨ぎグループの判定漏れ}：
    検出精度および人物バウンディングボックスの形成速度は改善されたものの、アプローチ09同様、物理的な横並び距離だけでは「スマートフォン操作や会話によって前後の列に跨って構成されているグループ（カップルや複数人のグループ）」を識別できない問題は未解決のまま残された。
\end{enumerate}
